%% title: Hierarchy: coarse parameters against coarse supervision
%% target: Sections/Results.tex
%% anchor: \resulttable{main_comparison_calibrated}
%% action: after
%% why: Splits the head comparison into its two questions. Per-level parameters never help in-distribution; per-level losses help genus and family, in-distribution on every backbone and under shift on two of three. At species under shift the ranking depends on the backbone.
%% source: /data/au761367/paper_numbers/13.md (campaign, seeds 42-44); journal/2026-10-23-which-recipe-does-the-merged-paper-rerun.md (discussion point 2, the per-rank table); RESULTS VS, VS2, VB0, VD
%% check: Rewritten 2026-11-09 with the rerun (owner's decision on C1). | The old trap claim (coarse supervision +2.9 on traps) reversed on V2-S; it is now stated per backbone. | Rename \subsection{Experiment 1} to \subsection{Hierarchical heads}\label{sec:results_heads}, which step 04 references. | AR: no rerun; the row is dropped; say in one clause that AR lost by 20 points in earlier runs if Asger wants it kept. | Trap ranges follow step 11's rule (non-overlapping three-seed ranges).
%% status: ready
\paragraph{Coarse parameters against coarse supervision.}
The heads differ in two ways that the comparison above confounds: whether genus and family have parameters of their own (I and TD), and whether they receive a loss of their own (every head except F).
F and BU read genus and family by summing species probabilities, so they differ only in the coarse loss.
\Cref{tab:heads_ranks} separates the two questions on three backbones, at every rank, in-distribution and on the trap images.

\begin{table}[htb]
    \centering
    \caption{Heads on three backbones, mean of three seeds. Test fold: macro-$\text{F}_1$ at zero threshold. Traps: macro-$\text{F}_1$ on all 47,905 trap crops, calibrated on 10\,\% of the (trap, night) groups. Bold: best on that backbone and column.}
    \label{tab:heads_ranks}
    \begin{tabular}{llcccccc}
        \toprule
        & & \multicolumn{3}{c}{test fold} & \multicolumn{3}{c}{traps} \\
        backbone & head & species & genus & family & species & genus & family \\
        \midrule
        EfficientNetV2-S & F & \textbf{0.889} & 0.949 & 0.969 & 0.354 & 0.468 & 0.522 \\
        & BU & 0.887 & \textbf{0.956} & \textbf{0.976} & 0.331 & 0.459 & 0.511 \\
        & I & 0.881 & 0.947 & 0.965 & \textbf{0.378} & \textbf{0.484} & \textbf{0.534} \\
        & TD & 0.842 & 0.930 & 0.960 & 0.327 & 0.433 & 0.525 \\
        \midrule
        EfficientNet-B0 & F & \textbf{0.875} & \textbf{0.940} & 0.961 & 0.280 & 0.391 & 0.423 \\
        & BU & 0.864 & \textbf{0.940} & \textbf{0.965} & \textbf{0.332} & \textbf{0.447} & \textbf{0.496} \\
        & I & 0.854 & 0.923 & 0.940 & 0.326 & 0.433 & 0.455 \\
        & TD & 0.788 & 0.887 & 0.934 & 0.267 & 0.360 & 0.482 \\
        \midrule
        DINOv3 ConvNeXt-L & F & \textbf{0.920} & 0.966 & 0.981 & 0.405 & 0.516 & 0.544 \\
        & BU & 0.918 & \textbf{0.968} & \textbf{0.984} & 0.416 & \textbf{0.548} & \textbf{0.587} \\
        & I & 0.917 & 0.966 & 0.981 & \textbf{0.424} & 0.534 & 0.571 \\
        & TD & 0.906 & 0.963 & 0.980 & 0.402 & 0.528 & 0.572 \\
        \bottomrule
    \end{tabular}
\end{table}

Coarse parameters do not help in-distribution: on every backbone F or BU is best at every rank of the test fold, and TD is lowest at every rank.
Coarse supervision raises family on the test fold on every backbone (BU over F by 0.3 to 0.7 points) and genus on two (0.2 and 0.7; unchanged on EfficientNet-B0), at a small cost at species (0.2 to 1.1 points, larger than the three-seed spread of at most 0.14).
On the trap images it raises genus and family by 5.6 and 7.3 points on EfficientNet-B0 and by 3.2 and 4.3 on DINOv3, beyond the three-seed spread; on EfficientNetV2-S the difference is within it.
At species under shift no head wins everywhere: coarse supervision is worth 5.2 points on EfficientNet-B0, 1.1 (within the seed spread) on DINOv3 and $-2.3$ on EfficientNetV2-S, while the independent head is best or within the spread of the best on all three.
Summing species probabilities therefore gives the best coarse predictions in-distribution, and supervising those sums gives better genus and family predictions on two of three backbones under shift: the predictions that back-off and abstention (\cref{sec:results_backoff}) fall back on.
